% !TEX root = ../memoria.tex
\chapter{Implementación}
\label{cap:capitulo4}

\section{Implementación de la Simulación en MATLAB}
\label{sec:impl_matlab}

\subsection{Visión general del motor de simulación}
\label{sec:matlab_vision}

La simulación radioeléctrica dispone de dos motores con funciones diferentes. El script
MATLAB \texttt{simulacion\_cobertura\_5g.m}, junto con
\texttt{exportar\_json.m}, genera los escenarios de referencia con Longley--Rice o con
el \textit{fallback} disponible en MATLAB. Para las nuevas peticiones web, el backend
Spring Boot ejecuta un motor Java FSPL+Deygout y guarda también la salida en JSON. Los
resultados de ambos motores comparten el propósito de visualizar RSRP, pero no son
intercambiables ni utilizan necesariamente el mismo modelo o terreno.

Decidimos dividir el script en nueve bloques ejecutables de forma independiente porque
durante el desarrollo fue clave poder validar cada paso por separado antes de encadenarlos.
El Cuadro~\ref{tab:bloques_matlab} los resume.

\begin{table}[H]
\centering
\caption{Estructura de bloques del script \texttt{simulacion\_cobertura\_5g.m}.}
\label{tab:bloques_matlab}
\small
\begin{tabularx}{\textwidth}{|c|l|X|}
\hline
\textbf{Bloque} & \textbf{Nombre} & \textbf{Función} \\
\hline
0 & Limpieza del entorno        & \texttt{clear}, \texttt{clc}, \texttt{close all} \\
1 & Parámetros de configuración & Define antena, frecuencia, potencia y cuadrícula \\
2 & Cuadrícula de receptores    & Genera el \textit{meshgrid} de puntos de muestreo \\
3 & Elevación SRTM              & Carga e interpola el MDT en la cuadrícula \\
4 & Modelo de propagación       & Calcula $P_{rx}$ en cada punto (LR-ITM o COST-231) \\
5 & Cálculo del RSRP            & Convierte $P_{rx}$ a RSRP según 3GPP TS 38.215 \\
6 & Visualización               & Genera \textit{heatmaps} y estadísticas de cobertura \\
7 & Exportación a JSON          & Serializa los resultados para el backend \\
8 & Prueba de referencia FSPL  & Comprueba casos teóricos bajo hipótesis de espacio libre \\
\hline
\end{tabularx}
\end{table}

El flujo es lineal: parámetros $\rightarrow$ cuadrícula $\rightarrow$ elevación
$\rightarrow$ $P_{rx}$ $\rightarrow$ RSRP $\rightarrow$ JSON. El fichero de salida, nombrado
con un UUID v4, es el único artefacto que cruza la frontera hacia el resto del sistema.

\subsection{Parámetros del sistema radiante (Bloque 1)}
\label{sec:matlab_params}

Todos los parámetros configurables están agrupados al inicio del script para que sea fácil
lanzar distintos escenarios sin buscar constantes dispersas por el código. El escenario de
referencia modela un \textit{gNodeB} típico de la banda FR1 n78 (3,4--3,8~GHz), que es
la banda de despliegue masivo de 5G en Europa, emplazado en el edificio de Rectorado del
campus URJC de Fuenlabrada. Los valores del Cuadro~\ref{tab:params_antena} son los parámetros adoptados para este
escenario de referencia; no se presentan como una instalación desplegada y medida.

\begin{table}[H]
\centering
\caption{Parámetros de la estación base (escenario URJC Fuenlabrada).}
\label{tab:params_antena}
\begin{tabular}{|l|c|c|}
\hline
\textbf{Parámetro} & \textbf{Símbolo} & \textbf{Valor} \\
\hline
Frecuencia central           & $f_c$         & 3,5~GHz (banda n78) \\
Potencia de transmisión      & $P_{tx}$      & 40~dBm (10~W) \\
Ganancia antena transmisora  & $G_{tx}$      & 15~dBi (sectorial) \\
Ganancia antena receptora    & $G_{rx}$      & 0~dBi (omnidireccional) \\
Pérdidas de cable            & $L_{cable}$   & 2~dB \\
Altura antena transmisora    & $h_{Tx}$      & 30~m \\
Altura antena receptora      & $h_{Rx}$      & 1,5~m \\
Emplazamiento (latitud)      & $\phi_0$      & 40,2897°~N \\
Emplazamiento (longitud)     & $\lambda_0$   & 3,8244°~W \\
\hline
\end{tabular}
\end{table}

Con estos valores, la EIRP queda en:

\begin{equation}
\text{EIRP} = P_{tx} + G_{tx} - L_{cable} = 40 + 15 - 2 = 53~\text{dBm}
\label{eq:eirp}
\end{equation}

En el escenario MATLAB de referencia, el área tiene 3~km de lado y se generan
$60 \times 60 = 3600$ puntos. Al incluir ambos extremos, el paso efectivo es
$3000/(60-1)\approx 50{,}85$~m. El generador actual persiste ese valor efectivo en el JSON;
los ficheros históricos ya existentes no se regeneraron. Esta configuración es una
elección de muestreo; no se atribuye un tiempo de ejecución sin una medición reproducible.

\begin{lstlisting}[language=MATLAB, caption={Parámetros de configuración de la simulación (Bloque 1).}, label={cod:bloque1}]
fc_Hz       = 3.5e9;        % Frecuencia central [Hz] - banda n78
Ptx_dBm     = 40;           % Potencia de transmisión [dBm] = 10 W
Ptx_W       = 10^((Ptx_dBm - 30) / 10);
Gtx_dBi     = 15;           % Ganancia antena Tx [dBi]
Grx_dBi     = 0;            % Ganancia antena Rx [dBi]
hTx_m       = 30;           % Altura antena Tx [m]
hRx_m       = 1.5;          % Altura antena Rx [m]
L_cable_dB  = 2;            % Pérdidas de cable [dB]
EIRP_dBm    = Ptx_dBm + Gtx_dBi - L_cable_dB;

lat0 = 40.2897;   lon0 = -3.8244;   % Campus URJC Fuenlabrada
areaKm        = 3;       % Lado del área cuadrada [km]
nPuntos       = 60;      % Muestras por eje, incluidos ambos extremos
gridSpacing_m = areaKm * 1000 / (nPuntos - 1);  % paso efectivo: aprox. 50.85 m
\end{lstlisting}

\subsection{Generación de la cuadrícula de receptores (Bloque 2)}
\label{sec:matlab_grid}

Para situar los 3600 receptores sobre el mapa hay que convertir el área en kilómetros a
incrementos en grados geográficos. Usamos la aproximación lineal estándar para latitudes
medias:

\begin{equation}
\Delta\phi = \frac{a/2}{111}~\text{grados}, \qquad
\Delta\lambda = \frac{a/2}{111 \cdot \cos\phi_0}~\text{grados}
\label{eq:deltas}
\end{equation}

donde $a = 3$~km. La conversión es una aproximación local adecuada para delimitar el área de este escenario;
no sustituye una proyección cartográfica ni se le atribuye una cota de error medida.

Con \texttt{linspace} construimos los vectores de coordenadas y con \texttt{meshgrid}
la cuadrícula 2D. Después aplanamos las matrices \texttt{LATS} y \texttt{LONS} a vectores
columna de 3600 elementos, que entran directamente en los bloques de elevación y propagación.

\begin{lstlisting}[language=MATLAB, caption={Generación de la cuadrícula de receptores (Bloque 2).}, label={cod:bloque2}]
deltaLat_deg = (areaKm / 2) / 111;
deltaLon_deg = (areaKm / 2) / (111 * cosd(lat0));

lats_vec = linspace(lat0 - deltaLat_deg, lat0 + deltaLat_deg, nPuntos);
lons_vec = linspace(lon0 - deltaLon_deg, lon0 + deltaLon_deg, nPuntos);

[LONS, LATS] = meshgrid(lons_vec, lats_vec);
lats_flat = LATS(:);   % Vectorizar: 3600 x 1
lons_flat = LONS(:);
\end{lstlisting}

\subsection{Datos de elevación del terreno (Bloque 3)}
\label{sec:matlab_srtm}

Los escenarios MATLAB existentes declaran \texttt{SRTM\_90m\_global} en sus metadatos.
Aunque el script puede leer un GeoTIFF local y trata de configurarlo en
\texttt{siteviewer}, esa posibilidad no demuestra que se utilizara en las salidas
almacenadas. El backend Java lee \texttt{matlab/data/srtm\_fuenlabrada.tif}; sus tags
GeoTIFF indican WGS~84, una escala angular de un segundo de arco
($0{,}00027778^\circ$ por píxel) y un raster de $3601\times3601$ muestras. Esta resolución
nominal equivale aproximadamente a 31~m en latitud y varía en longitud. El nombre del
fichero sugiere SRTM, pero
no se conserva una ficha de descarga, proveedor comprobable ni suma de verificación para
establecer su procedencia documental. Las pruebas confirman que el fichero se puede leer,
que su extensión se expone y que se rechazan coordenadas exteriores.

Para interpolar la elevación en cada punto de la cuadrícula usamos \texttt{readgeoraster}
y \texttt{geointerp} del \textit{Mapping Toolbox}:

\begin{lstlisting}[language=MATLAB, caption={Carga e interpolación del MDT SRTM (Bloque 3).}, label={cod:srtm}]
[elev_data, R] = readgeoraster(srtm_file);
elev_data = double(elev_data);
elev_data(elev_data < -100) = 0;          % Eliminar valores sin dato
elev_grid = geointerp(elev_data, R, lats_flat, lons_flat, 'linear');
\end{lstlisting}

Si el fichero no está disponible, el script avisa con un \textit{warning} y sustituye
\texttt{elev\_grid} por ceros, asumiendo terreno plano. Esto resultó muy útil en las
primeras fases del desarrollo, cuando todavía no habíamos descargado el tile. La variable
\texttt{tiene\_srtm} registra cuál de los dos caminos se tomó y queda reflejada en el campo
\texttt{"terrain\_data"} del JSON de salida.

\subsection{Modelo de propagación radioeléctrica (Bloque 4)}
\label{sec:matlab_propagacion}

Este es el bloque central: calcula $P_{rx}$ en los 3600 puntos aplicando un modelo de
propagación semiempírico. Al arrancar, el script comprueba qué \textit{toolboxes} están
instalados y selecciona el mejor disponible automáticamente.

\subsubsection{Modelo Longley-Rice (ITM)}
\label{sec:matlab_lr}

Elegimos el \textit{Irregular Terrain Model} (ITM) o modelo Longley-Rice
\cite{longley1968transmission} como motor de referencia de los escenarios MATLAB por su
incorporación del perfil del terreno en el cálculo de pérdidas. A diferencia
de los modelos puramente empíricos, evalúa la difracción sobre las irregularidades del
terreno (parametrizada por $\Delta h$), la reflexión en el suelo y la curvatura de la
Tierra para cada enlace TX-RX individual. Para redes desplegables en entornos con orografía
variable, ignorar eso sería un error de bulto.

El modelo, desarrollado por A.~G.~Longley e Y.~T.~Rice en 1968 en el ITS de la NTIA, lleva
décadas siendo la referencia en planificación de redes de radiodifusión y telecomunicaciones
militares. En MATLAB está disponible desde R2019b a través del \textit{Communications Toolbox}:

\begin{lstlisting}[language=MATLAB, caption={Cálculo con Longley-Rice ITM (Bloque 4).}, label={cod:lr}]
tx = txsite('Name', 'gNodeB_URJC', ...
    'Latitude',             lat0,  'Longitude',       lon0, ...
    'AntennaHeight',        hTx_m, 'TransmitterFrequency', fc_Hz, ...
    'TransmitterPower',     Ptx_W);

pm = propagationModel('longley-rice');

rx_sites = rxsite('Latitude',      lats_flat, ...
                  'Longitude',     lons_flat, ...
                  'AntennaHeight', hRx_m);

Prx_dBm = sigstrength(rx_sites, tx, pm);
\end{lstlisting}

\texttt{sigstrength} devuelve $P_{rx}$ en dBm para cada receptor. El script intenta
configurar un terreno local mediante las interfaces de \texttt{rfprop.TerrainModel} o
\texttt{siteviewer}; si fallan, utiliza el terreno global de MATLAB y registra
\texttt{SRTM\_90m\_global}. Los tres JSON de referencia disponibles declaran esta última
fuente. En consecuencia, no se atribuye a esas simulaciones el uso efectivo del GeoTIFF
local ni una mejora de precisión derivada de una resolución de 30~m.

\subsubsection{COST-231 Hata como modelo de respaldo}
\label{sec:matlab_cost231}

Cuando el \textit{Communications Toolbox} no está instalado, el script recurre a COST-231
Hata extendido~\cite{cost231finalreport1999}. Este modelo solo necesita funciones
matemáticas básicas de MATLAB. La fórmula de
pérdidas para entorno urbano medio ($C_m = 3$~dB) es:

\begin{equation}
\begin{split}
L_{50}~[\text{dB}] ={}& 46{,}3 + 33{,}9 \log_{10}f_c - 13{,}82 \log_{10}h_{Tx} \\
                      &- a(h_{Rx}) + (44{,}9 - 6{,}55 \log_{10}h_{Tx})\log_{10}d + C_m
\end{split}
\label{eq:cost231}
\end{equation}

con $f_c$ en MHz, $d$ en km y $a(h_{Rx})$ el factor de corrección de la antena receptora:

\begin{equation}
a(h_{Rx}) = (1{,}1 \log_{10}f_c - 0{,}7)\,h_{Rx} - (1{,}56 \log_{10}f_c - 0{,}8)
\label{eq:a_hrx}
\end{equation}

De ahí, $P_{rx} = \text{EIRP} + G_{rx} - L_{50}$.

El problema es que COST-231 Hata fue desarrollado para frecuencias de hasta 2000~MHz.
Su aplicación a 3,5~GHz queda fuera de ese rango, por lo que el resultado de ese
\textit{fallback} no se presenta como predicción cuantitativamente validada.

\subsubsection{FSPL como referencia teórica}
\label{sec:matlab_fspl}

Sea cual sea el modelo activo, también calculamos FSPL (\textit{Free Space Path Loss})
como referencia teórica de propagación en espacio libre:

\begin{equation}
\text{FSPL}~[\text{dB}] = 20\log_{10}d + 20\log_{10}f_c + 20\log_{10}\!\left(\frac{4\pi}{c}\right)
\label{eq:fspl}
\end{equation}

donde $d$ está en metros y $f_c$ en Hz. La diferencia respecto al modelo activo refleja
las hipótesis adicionales de este último y no aísla por sí sola el efecto del terreno.
Esta comparación se calcula en el bloque 8.

\subsection{Cálculo del RSRP (Bloque 5)}
\label{sec:matlab_rsrp}

$P_{rx}$ representa la potencia total integrada sobre todo el ancho de banda del canal.
Para SS-RSRP, 3GPP TS~38.215 define la media lineal de las contribuciones de potencia de
los \textit{Resource Elements} que transportan el \textit{Secondary Synchronization Signal}
dentro del ancho de banda de medida \cite{3gpp_ts38215_2022}. El cálculo de este prototipo
no reproduce una medición de terminal: estima RSRP a partir de $P_{rx}$ suponiendo una
distribución uniforme de potencia entre los RE considerados, mediante la normalización
siguiente:

\begin{equation}
\text{RSRP}~[\text{dBm}] = P_{rx}~[\text{dBm}] - 10\log_{10}(N_{RB} \cdot 12)
\label{eq:rsrp}
\end{equation}

donde el factor 12 es el número de subportadoras por \textit{Resource Block}. El generador
MATLAB conserva factores de normalización heredados: usa 66 para frecuencias de al menos
3~GHz y 52 por debajo. Como los escenarios no especifican ancho de banda ni SCS, estos
factores no se interpretan como recursos asignados ni como una configuración normalizada
de 100~MHz/30~kHz. Para el factor 66, la corrección es:

\begin{equation}
\Delta_{\text{RSRP}} = 10\log_{10}(66 \times 12) = 10\log_{10}(792) \approx 28{,}99~\text{dB}
\label{eq:delta_rsrp}
\end{equation}

Los valores estimados se limitan al rango $[-140, -40]$~dBm y se reconstruyen en
matriz $60 \times 60$ para visualización y exportación. TS~38.133, apartado~10.1.6, especifica
un mapeo de mediciones RSRP a valores de informe \cite{etsi_ts138133_v153400_2026}; los
intervalos del Cuadro~\ref{tab:rsrp_umbrales} son, en cambio, umbrales visuales elegidos
por el prototipo. No son las categorías de ese mapeo ni permiten inferir por sí solos
calidad universal o garantía de servicio.

\begin{lstlisting}[language=MATLAB, caption={Cálculo del RSRP a partir de $P_{rx}$ (Bloque 5).}, label={cod:bloque5}]
if fc_Hz >= 3e9
    NRB_ASSUMED = 66;
else
    NRB_ASSUMED = 52;
end
delta_RSRP_dB = 10 * log10(NRB_ASSUMED * 12);

RSRP_dBm = Prx_dBm - delta_RSRP_dB;

RSRP_min = -140;  RSRP_max = -40;           % Límites de salida configurados
RSRP_dBm = max(min(RSRP_dBm, RSRP_max), RSRP_min);

RSRP_grid = reshape(RSRP_dBm, nPuntos, nPuntos);
\end{lstlisting}

\begin{table}[H]
\centering
\caption{Intervalos visuales de RSRP adoptados por el prototipo.}
\label{tab:rsrp_umbrales}
\begin{tabular}{|l|c|l|}
\hline
\textbf{Etiqueta visual} & \textbf{Rango RSRP} & \textbf{Descripción orientativa} \\
\hline
Muy alto      & $\geq -80$~dBm          & Extremo superior de la escala visual \\
Alto          & $[-90, -80)$~dBm        & Nivel superior respecto a los intervalos siguientes \\
Intermedio    & $[-100, -90)$~dBm       & Nivel intermedio de la escala visual \\
Bajo          & $[-110, -100)$~dBm      & Nivel inferior de la escala visual \\
Muy bajo      & $< -110$~dBm            & Extremo inferior de la escala visual \\
\hline
\end{tabular}
\end{table}

\subsection{Exportación de resultados a JSON (Bloque 7)}
\label{sec:matlab_json}

Con el RSRP calculado, el bloque 7 llama a \texttt{exportar\_json.m}. La función genera
un UUID v4 vía la JVM integrada en MATLAB (\texttt{java.util.UUID.randomUUID()}), empaqueta
los metadatos y los puntos de cobertura en una estructura, serializa todo con
\texttt{jsonencode} y escribe el fichero en \texttt{matlab/output/<uuid>.json}. No hay
dependencias externas: la JVM viene de serie con cualquier instalación de MATLAB.

El formato de salida lo diseñamos para que el backend Spring Boot pueda leerlo sin
transformaciones intermedias:

\begin{lstlisting}[caption={Estructura del fichero JSON generado por \texttt{exportar\_json.m}.}, label={cod:json}]
{
  "id": "a3f7c2b1-...",
  "name": "URJC_Fuenlabrada_3.5GHz_40dBm",
  "createdAt": "2026-04-15T12:00:00Z",
  "metadata": {
    "frequency_ghz": 3.5,
    "tx_power_dbm": 40,
    "antenna_height_m": 30,
    "tx_location": { "latitude": 40.2897, "longitude": -3.8244 },
    "model": "Longley-Rice",
    "terrain_data": "SRTM_90m_global"
  },
  "coveragePoints": [
    { "lat": 40.2764, "lng": -3.8356, "rsrp": -97.4 },
    ...
  ]
}
\end{lstlisting}

El array \texttt{coveragePoints} solo incluye los puntos con RSRP superior a $-130$~dBm.
El motor Java aplica el mismo umbral de exportación. Por tanto, el número de puntos
serializados puede ser inferior al total de muestras de la cuadrícula.

El Listado~\ref{cod:exportar_json} muestra el núcleo de la función
\texttt{exportar\_json.m}: la generación del UUID, la construcción de la estructura,
la marca de tiempo ISO~8601 en UTC y la serialización en la carpeta
\texttt{matlab/output/}.

\begin{lstlisting}[language=MATLAB, caption={Fragmento principal de \texttt{exportar\_json.m}.}, label={cod:exportar_json}]
function fichero_json = exportar_json(nombre, lats, lons, rsrp_dBm, params)
    rsrp_dBm = max(min(double(rsrp_dBm), -40), -140);  % Clamp
    uuid     = char(java.util.UUID.randomUUID().toString());
    ahora    = datetime('now', 'TimeZone', 'UTC');
    createdAt = char(ahora, "yyyy-MM-dd'T'HH:mm:ss'Z'");

    metadata.frequency_ghz    = double(params.frecuencia_GHz);
    metadata.tx_power_dbm     = double(params.potencia_tx_dBm);
    metadata.antenna_height_m = double(params.altura_tx_m);
    metadata.tx_location      = struct('latitude',  params.lat_emplazamiento, ...
                                       'longitude', params.lon_emplazamiento);
    metadata.model            = char(string(params.modelo_propagacion));
    metadata.terrain_data     = char(string(params.datos_elevacion));

    idx_validos = rsrp_dBm >= -130;  % Filtrar puntos sin señal útil
    coveragePoints = struct('lat', num2cell(lats(idx_validos)), ...
                            'lng', num2cell(lons(idx_validos)), ...
                            'rsrp', num2cell(rsrp_dBm(idx_validos)));

    sim_struct = struct('id', uuid, 'name', nombre, ...
        'createdAt', createdAt, 'metadata', metadata, ...
        'coveragePoints', coveragePoints);

    json_str = jsonencode(sim_struct, 'PrettyPrint', true);
    fichero_json = fullfile(output_dir, [uuid, '.json']);
    fid = fopen(fichero_json, 'w', 'n', 'UTF-8');
    fprintf(fid, '%s', json_str);  fclose(fid);
end
\end{lstlisting}

\section{Backend: API REST con Spring Boot}
\label{sec:impl_backend}

El backend expone los ficheros JSON existentes y genera nuevas simulaciones bajo demanda
para que el frontend pueda consumirlas. La API REST se implementa con Spring Boot~3.2.4 y
Java~17, empaquetada como un JAR ejecutable que no requiere servidor de aplicaciones
externo. Las peticiones de consulta no mantienen estado de sesión; la creación, en cambio,
es una operación \texttt{POST} que persiste una nueva salida JSON en el directorio
configurado.

\subsection{Estructura del proyecto}
\label{sec:backend_estructura}

El proyecto sigue la estructura canónica de Maven con Spring Boot. El fichero
\texttt{pom.xml} declara tres dependencias principales: \texttt{spring-boot-starter-web}
(servidor embebido Tomcat, Jackson para JSON y el framework web de Spring MVC),
\texttt{springdoc-openapi-starter-webmvc-ui}~2.3.0 (generación automática de la
documentación OpenAPI~3 y la interfaz Swagger~UI) y \texttt{spring-boot-starter-test}
(JUnit~5, Mockito y utilidades de test de Spring).

El código Java se organiza en cinco paquetes siguiendo el patrón clásico de capas:

\begin{table}[H]
\centering
\caption{Estructura de paquetes del backend.}
\label{tab:backend_paquetes}
\begin{tabularx}{\textwidth}{|l|X|}
\hline
\textbf{Paquete} & \textbf{Responsabilidad} \\
\hline
\texttt{controller} & Controlador REST: define endpoints y formato de respuesta \\
\texttt{service}    & Lógica de negocio: lectura y validación de ficheros JSON \\
\texttt{model}      & DTOs: \texttt{SimulationSummaryDto} (record Java) \\
\texttt{exception}  & Excepciones de dominio y manejador global \\
\texttt{config}     & Configuración CORS para desarrollo local \\
\hline
\end{tabularx}
\end{table}

La clase de arranque \texttt{Gemelo5gApplication} es un \texttt{@SpringBootApplication}
estándar que delega en \texttt{SpringApplication.run()} el escaneo de componentes y la
configuración automática del contenedor.

El \texttt{SidebarComponent} muestra el histórico de simulaciones y ofrece una acción de
borrado con confirmación explícita. Tras confirmarla, el frontend envía
\texttt{DELETE /api/simulations/\{id\}}; el backend comprueba la identidad entre el UUID
solicitado y el contenido del JSON antes de eliminarlo, y responde con conflicto si no
coinciden. En caso de éxito, el frontend retira la entrada, conserva la selección activa
si sigue disponible o selecciona otra, y limpia el punto seleccionado. En comparación,
\texttt{ComparisonService} reconcilia B tras el borrado, impide que una respuesta tardía
de una carga anterior restaure una selección obsoleta y no permite comparar A consigo
misma. La colocación de un transmisor queda desactivada durante la comparación. Estas
transiciones están cubiertas por pruebas automatizadas y por la verificación funcional
local en navegador descrita en el Capítulo~\ref{cap:capitulo5}; no se dispone de una suite
de navegador integrada en CI.

\subsection{Servicio de simulaciones}
\label{sec:backend_service}

\texttt{SimulationService} es el componente central del backend. Anotado con
\texttt{@Service}, recibe por inyección de dependencias la ruta del directorio de datos
(propiedad \texttt{simulations.data-dir} de \texttt{application.properties}) y un
\texttt{ObjectMapper} de Jackson para la deserialización JSON.

El directorio de datos se normaliza en el constructor con
\texttt{Paths.get(path).toAbsolutePath().normalize()} para prevenir ataques de
\textit{path traversal}: cualquier ruta resuelta que no empiece por el directorio base
se rechaza con un \texttt{SimulationNotFoundException}. Además, antes de acceder al
sistema de ficheros, el método \texttt{getById} valida que el identificador recibido
sea un UUID~v4 conforme a RFC~4122 mediante la expresión regular:

\begin{center}
\texttt{[0-9a-fA-F]\{8\}-[0-9a-fA-F]\{4\}-}\\
\texttt{[0-9a-fA-F]\{4\}-[0-9a-fA-F]\{4\}-[0-9a-fA-F]\{12\}}
\end{center}

Si el identificador no supera ninguna de las dos validaciones, se lanza la misma
excepción para no revelar información sobre la estructura del sistema de ficheros.

\begin{lstlisting}[language=Java, caption={Métodos principales de \texttt{SimulationService}.}, label={cod:service}]
@Service
public class SimulationService {
    private final Path dataDir;
    private final ObjectMapper mapper;

    public SimulationService(
            @Value("${simulations.data-dir}") String dataDirPath,
            ObjectMapper mapper) {
        this.dataDir = Paths.get(dataDirPath).toAbsolutePath().normalize();
        this.mapper  = mapper;
    }

    public List<SimulationSummaryDto> listAll() throws IOException {
        if (!Files.isDirectory(dataDir)) return List.of();
        try (Stream<Path> files = Files.list(dataDir)) {
            return files.filter(p -> p.toString().endsWith(".json"))
                        .map(this::parseSummary).toList();
        }
    }

    private static final String UUID_REGEX =
        "[0-9a-fA-F]{8}-[0-9a-fA-F]{4}-[0-9a-fA-F]{4}"
      + "-[0-9a-fA-F]{4}-[0-9a-fA-F]{12}";

    public JsonNode getById(String id) throws IOException {
        if (!id.matches(UUID_REGEX))
            throw new SimulationNotFoundException(id);
        Path file = dataDir.resolve(id + ".json").normalize();
        if (!file.startsWith(dataDir) || !Files.exists(file))
            throw new SimulationNotFoundException(id);
        return mapper.readTree(file.toFile());
    }
}
\end{lstlisting}

El método \texttt{listAll} recorre el directorio con \texttt{Files.list()}, filtra los
ficheros \texttt{.json} y mapea cada uno a un \texttt{SimulationSummaryDto} que contiene
solo los campos ligeros (\texttt{id}, \texttt{name}, \texttt{createdAt},
\texttt{metadata}), sin incluir el array \texttt{coveragePoints}. De este modo, el
listado no satura al cliente cuando hay múltiples simulaciones almacenadas.

\subsection{Motor Java para simulaciones bajo demanda}
\label{sec:backend_java_engine}

En la fase~1, la operación \texttt{POST /api/simulations} recibe un nombre, frecuencia, altura del
receptor, tamaño de área, paso de cuadrícula, pérdidas de cable y un transmisor. El modelo
\texttt{SimulationRequest} limita la frecuencia a 0,4--7,125~GHz, el área a
0,1--50~km, el paso solicitado a 10--10\,000~m y la cuadrícula a 100\,000 puntos como
máximo, además de un límite estimado de 20 millones de muestras de perfiles.
La rejilla incluye ambos extremos: calcula
$\lceil\text{área}/\text{paso solicitado}\rceil+1$ puntos por eje y conserva en los
metadatos los pasos solicitado y efectivo. El alcance documentado es un único transmisor.

Para cada punto de la cuadrícula, \texttt{PropagationEngine} calcula la pérdida en espacio
libre y añade la pérdida de difracción de Deygout sobre un perfil de elevación muestreado
del GeoTIFF local~\cite{deygout1966multiple}. La curvatura terrestre se aproxima mediante
un radio efectivo de $4/3$ del radio medio terrestre. A partir de la EIRP se obtiene un
proxy de potencia, etiquetado RSRP en la interfaz. Este valor se limita al intervalo
$[-140,-40]$~dBm y se omiten del JSON los puntos por debajo de
$-130$~dBm. Esta implementación es distinta de Longley--Rice y no modela
edificios ni reflexiones en fachadas. El proxy no es una medida física de SS-RSRP.
Las extensiones posteriores integradas localmente y sus límites se describen en la
Sección~\ref{sec:avances_posteriores}.

\texttt{SimulationCreationService} crea un UUID y una marca temporal, incorpora los
parámetros y metadatos del motor (\texttt{FSPL+Deygout} y origen del terreno)
y persiste el JSON en el directorio configurado. El valor predeterminado de
\texttt{terrain.source} es \texttt{Raster local (procedencia no especificada)};
la configuración de Compose especifica terreno sintético para su raster de demo.
La etiqueta no verifica por sí sola la procedencia del fichero. Una versión anterior
atribuyó indebidamente \texttt{SRTM\_30m\_local} a terreno sintético. La escritura
se realiza primero en un fichero temporal y después se mueve al nombre definitivo, usando
un movimiento atómico cuando el sistema de archivos lo permite. El resultado se devuelve
al cliente para que la interfaz lo seleccione y visualice.

\subsection{Apoyos de planificación y opciones radio}
\label{sec:apoyos_planificacion}

Además del recorrido principal, una implementación local permite importar observaciones
compatibles con el proxy de potencia y calibrar un desplazamiento, conservando el
escenario original. Se separan las observaciones de ajuste y reserva, se emparejan con
los puntos cercanos de la malla y se informa el sesgo y el RMSE antes y después del
ajuste. Los ensayos descritos en el capítulo~\ref{cap:capitulo5} son sintéticos: no
convierten el proxy en SS-RSRP físico ni acreditan mejora en campo.

La recomendación de emplazamiento evalúa candidatos y objetivos discretos sobre un
área cuadrada o un polígono simple. La selección voraz favorece al candidato que cubre
más objetivos pendientes; depende de la malla y del umbral y no garantiza un óptimo
global ni cobertura continua. El resultado es una ayuda para el operador, no una
sustitución de las restricciones y la inspección del lugar.

Como complementos se han ensayado un patrón direccional aproximado configurable por
azimut e inclinación, una pérdida de \textit{clutter} introducida manualmente y un
perfil de enlace con primera zona de Fresnel. No representan, respectivamente, un
diagrama medido de antena ni un mapa real de usos del suelo; el perfil depende de la
procedencia y calidad del raster. La integración con el frontend y backend reorganizados
debe comprobarse de nuevo.

\subsection{Controlador REST}
\label{sec:backend_controller}

\texttt{SimulationController} expone operaciones de consulta y de creación bajo el prefijo
\texttt{/api/simulations}. El Cuadro~\ref{tab:impl_endpoints} las resume.

\begin{table}[H]
\centering
\caption{Endpoints principales de la API REST.}
\label{tab:impl_endpoints}
\small
\begin{tabularx}{\textwidth}{|c|>{\raggedright\arraybackslash}p{5.7cm}|X|}
\hline
\textbf{Método} & \textbf{URL} & \textbf{Descripción} \\
\hline
POST & \texttt{/api/simulations} & Valida parámetros, calcula una simulación Java y persiste el JSON \\
DELETE & \texttt{/api/simulations/\{id\}} & Elimina el JSON tras confirmar; rechaza UUID inválido o identidad inconsistente \\
GET & \texttt{/api/simulations} & Listado de simulaciones (sin puntos) \\
GET & \texttt{/api/simulations/\{id\}} & JSON completo con \texttt{coveragePoints} \\
GET & \texttt{/api/simulations/\{id\}/export} & Descarga como fichero adjunto \\
\hline
\end{tabularx}
\end{table}

El endpoint de listado soporta paginación opcional: si el cliente envía los parámetros
\texttt{?page=0\&size=10}, el controlador calcula los índices de corte sobre la lista
completa y devuelve un DTO \texttt{PageResponse} con los campos \texttt{content},
\texttt{totalElements}, \texttt{page}, \texttt{size} y \texttt{totalPages}. Si no se
envía el parámetro \texttt{page}, devuelve la lista completa para mantener la
compatibilidad con clientes simples.

\begin{lstlisting}[language=Java, caption={Fragmento de endpoints del controlador REST.}, label={cod:controller}]
@RestController
@RequestMapping("/api/simulations")
public class SimulationController {
    private final SimulationService service;
    private final SimulationCreationService creationService;

    @PostMapping(consumes = MediaType.APPLICATION_JSON_VALUE,
                 produces = MediaType.APPLICATION_JSON_VALUE)
    public JsonNode createSimulation(
            @Valid @RequestBody SimulationRequest request) throws IOException {
        return creationService.create(request);
    }

    @DeleteMapping("/{id}")
    public ResponseEntity<Void> deleteSimulation(@PathVariable String id) {
        service.deleteById(id);
        return ResponseEntity.noContent().build();
    }

    @GetMapping(produces = MediaType.APPLICATION_JSON_VALUE)
    public ResponseEntity<?> listSimulations(
            @RequestParam(required = false) Integer page,
            @RequestParam(required = false, defaultValue = "10") int size)
            throws IOException {
        List<SimulationSummaryDto> all = service.listAll();
        if (page == null) return ResponseEntity.ok(all);
        int total = all.size();
        int from  = Math.min(page * size, total);
        int to    = Math.min(from + size, total);
        return ResponseEntity.ok(
            new PageResponse<>(all.subList(from, to), total, page, size,
                (int) Math.ceil((double) total / size)));
    }

    @GetMapping(value = "/{id}", produces = MediaType.APPLICATION_JSON_VALUE)
    public JsonNode getSimulation(@PathVariable String id) throws IOException {
        return service.getById(id);
    }

    @GetMapping(value = "/{id}/export", produces = MediaType.APPLICATION_JSON_VALUE)
    public ResponseEntity<Resource> exportSimulation(@PathVariable String id)
            throws IOException {
        Resource resource = service.getFileResource(id);
        return ResponseEntity.ok()
            .contentType(MediaType.APPLICATION_JSON)
            .header(HttpHeaders.CONTENT_DISPOSITION,
                    "attachment; filename=\"" + id + ".json\"")
            .body(resource);
    }
}
\end{lstlisting}

El endpoint de exportación añade la cabecera HTTP de descarga para que el navegador
descargue el JSON como fichero en vez de mostrarlo en línea.
Internamente delega en \texttt{SimulationService.getFileResource()}, que devuelve un
\texttt{FileSystemResource} apuntando al fichero original.

\subsection{Gestión de errores}
\label{sec:backend_errores}

La gestión centralizada de errores se apoya en dos clases. Primero,
\texttt{SimulationNotFoundException} extiende \texttt{RuntimeException} y lleva la
anotación \texttt{@ResponseStatus(HttpStatus.NOT\_FOUND)}: cuando se lanza, Spring
devuelve automáticamente un código 404. Segundo, \texttt{GlobalExceptionHandler}
(anotado con \texttt{@RestControllerAdvice}) captura las excepciones no manejadas y
las transforma en respuestas JSON consistentes con estructura
\texttt{\{"error": "...", "detail": "..."\}}:

\begin{lstlisting}[language=Java, caption={Manejador global de excepciones.}, label={cod:exception}]
@RestControllerAdvice
public class GlobalExceptionHandler {
    @ExceptionHandler(SimulationNotFoundException.class)
    @ResponseStatus(HttpStatus.NOT_FOUND)
    public Map<String, String> handleNotFound(SimulationNotFoundException ex) {
        return Map.of("error", "Simulacion no encontrada",
                      "detail", ex.getMessage());
    }

    @ExceptionHandler(IOException.class)
    @ResponseStatus(HttpStatus.INTERNAL_SERVER_ERROR)
    public Map<String, String> handleIoException(IOException ex) {
        return Map.of("error", "Error leyendo datos de simulacion",
                      "detail", ex.getMessage());
    }
}
\end{lstlisting}

Este diseño garantiza que el frontend siempre recibe un JSON interpretable, nunca
una traza de pila ni una página de error HTML por defecto de Spring.

\subsection{Configuración y CORS}
\label{sec:backend_config}

El fichero \texttt{application.properties} contiene tres grupos de configuración:

\begin{lstlisting}[caption={Configuración del backend (\texttt{application.properties}).}, label={cod:appprops}]
server.port=8080
simulations.data-dir=../matlab/output
springdoc.swagger-ui.path=/swagger-ui.html
springdoc.api-docs.path=/api-docs
\end{lstlisting}

La propiedad \texttt{simulations.data-dir} apunta por defecto al directorio
\texttt{matlab/output} relativo a la carpeta del backend, de modo que el servidor
lee directamente los JSON que genera MATLAB sin copia intermedia. En producción o en
entornos con rutas distintas, basta con cambiar esta propiedad o pasarla como variable
de entorno (\texttt{SIMULATIONS\_DATA\_DIR}).

La clase \texttt{CorsConfig} registra una política CORS que permite peticiones desde
\texttt{http://localhost:4200} (el servidor de desarrollo de Angular) con los métodos
\texttt{GET}, \texttt{POST} y \texttt{DELETE}. Spring gestiona las peticiones
\textit{preflight} \texttt{OPTIONS} para comprobar los métodos solicitados, sin incluir
\texttt{OPTIONS} en la lista de métodos de aplicación configurados. Esta configuración es imprescindible durante el
desarrollo local, ya que navegador y backend corren en puertos distintos y sin la
cabecera \texttt{Access-Control-Allow-Origin} el navegador bloquearía las peticiones.

Finalmente, \texttt{springdoc-openapi} genera automáticamente la especificación
OpenAPI~3 a partir de las anotaciones \texttt{@Operation} y \texttt{@Tag} del
controlador. La interfaz Swagger~UI queda accesible en
\texttt{http://localhost:8080/swagger-ui.html} y permite probar los endpoints
directamente desde el navegador sin necesidad de herramientas externas como Postman.

\subsection{Tests unitarios}
\label{sec:backend_tests}

El informe Maven de fase~1, fechado el 21/11/2026, contabilizó 47 pruebas: 46 ejecutadas y
una comparación Java--MATLAB optativa omitida, sin fallos ni errores. El Cuadro~
\ref{tab:tests_backend} resume los recuentos por clase.

\begin{table}[H]
\centering
\caption{Recuento por clase del informe Maven de fase~1 (21/11/2026).}
\label{tab:tests_backend}
\small
\begin{tabularx}{\textwidth}{|l|X|X|}
\hline
\textbf{Clase de test} & \textbf{Tests} & \textbf{Ámbito} \\
\hline
\texttt{SimulationControllerTest} & 15 & Endpoints, errores, borrado y CORS \\
\texttt{SimulationRequestValidationTest} & 4 & Límites, estructura y transmisor nulo \\
\texttt{GeoTiffElevationModelTest} & 5 & PixelIsArea sintético, PixelIsPoint, lectura y extensión \\
\texttt{PropagationEngineTest} & 5 & FSPL, difracción y cálculo de la cuadrícula \\
\texttt{SimulationCreationServiceTest} & 1 & Creación y persistencia JSON \\
\texttt{SimulationServiceTest} & 16 & Consulta, listado, borrado y exportación \\
\texttt{MatlabScenarioComparisonTest} & 1 & Comparación optativa Java--MATLAB (omitida en esta ejecución) \\
\hline
\end{tabularx}
\end{table}

El XML JaCoCo de fase~1 (\texttt{backend/target/site/jacoco/jacoco.xml}, generado el
21/11/2026 a las 21:22) registr\'o 409 de 453 l\'ineas cubiertas (90,3~\%) y 137 de 192
ramas (71,4~\%) en 26 clases de producci\'on. La cobertura de c\'odigo es una medida
estructural: no acredita la exactitud de la predicci\'on radioel\'ectrica ni sustituye una
validaci\'on con mediciones.

En la validación del candidato integrado, Maven registró 82 casos: 79 ejecutados y
superados y tres omitidos, sin fallos. Los omitidos corresponden al SRTM local, la
recomendación con raster local y la comparación histórica. Un control selectivo
adicional Java--MATLAB sintético fue satisfactorio con diferencias inferiores a
$10^{-8}$~dB. Los resultados detallados y sus límites se presentan en el
Capítulo~\ref{cap:capitulo5}.

\section{Frontend: Aplicación Angular}
\label{sec:impl_frontend}

La capa de presentación es una aplicación web monopágina (\textit{Single Page
Application}) construida con Angular~21 y TypeScript~5.9. Su función es consumir la
API REST del backend y representar los datos de cobertura como un \textit{heatmap}
interactivo sobre un mapa geográfico. A diferencia de versiones anteriores de Angular,
la aplicación utiliza exclusivamente componentes \textit{standalone} (sin
\texttt{NgModule}) y gestión de estado basada en \textit{signals}, siguiendo las
prácticas recomendadas en Angular moderno.

\subsection{Arquitectura de componentes}
\label{sec:frontend_componentes}

La aplicación está compuesta por ocho componentes, cada uno independiente y declarado
con \texttt{standalone:~true}. El Cuadro~\ref{tab:componentes_angular} los enumera junto
con su responsabilidad.

\begin{table}[H]
\centering
\caption{Componentes de la aplicación Angular.}
\label{tab:componentes_angular}
\begin{tabularx}{\textwidth}{|>{\raggedright\arraybackslash}p{5.8cm}|X|}
\hline
\textbf{Componente} & \textbf{Responsabilidad} \\
\hline
\texttt{HeaderComponent}           & Barra superior con título y toggle de tema \\
\texttt{SidebarComponent}          & Selector de simulación, histórico, filtro RSRP y estadísticas \\
\texttt{MapComponent}              & Mapa Leaflet con \textit{heatmap} y marcador TX \\
\texttt{InfoPanelComponent}        & Detalle del punto seleccionado al hacer clic \\
\texttt{StatsPanelComponent}       & Resumen estadístico de la simulación activa \\
\texttt{CoverageLegendComponent}   & Leyenda de colores del \textit{heatmap} \\
\texttt{SimulationConfigComponent} & Formulario de parámetros, colocación del transmisor y ejecución del POST \\
\texttt{ComparisonViewComponent} & Comparación de dos escenarios en mapas y paneles de parámetros \\
\hline
\end{tabularx}
\end{table}

\texttt{AppComponent} actúa como raíz: importa los componentes
\texttt{Header}, \texttt{Sidebar}, \texttt{Map} e \texttt{InfoPanel}, y en su
método \texttt{ngOnInit} invoca \texttt{SimulationService.loadSimulations()} para
cargar el listado inicial desde el backend. Si la petición falla, se muestra un
\textit{banner} de error con botón de reintento que vuelve a llamar al mismo método.

\subsection{Servicio de simulaciones}
\label{sec:frontend_service}

\texttt{SimulationService} centraliza toda la comunicación con la API y el estado
de la aplicación. Está registrado con \texttt{providedIn:~'root'}, lo que lo convierte
en un \textit{singleton} inyectable en cualquier componente.

El estado interno se gestiona con cinco \textit{signals} de escritura
(\texttt{WritableSignal}): \texttt{simulations} (lista cargada),
\texttt{selectedSimulationId} (UUID activo), \texttt{rsrpThreshold} (umbral del
filtro, inicializado a $-140$~dBm), \texttt{selectedPoint} (punto clicado en el
mapa), \texttt{loading} y \texttt{error}. A partir de ellos se derivan cuatro
\textit{computed signals} de solo lectura, entre los que destaca
\texttt{filteredCoveragePoints}:

\begin{lstlisting}[language=JavaScript, caption={Servicio de simulaciones (fragmento de \texttt{simulation.service.ts}).}, label={cod:sim_service}]
@Injectable({ providedIn: 'root' })
export class SimulationService {
  private readonly http = inject(HttpClient);
  private simulations        = signal<Simulation[]>([]);
  private selectedSimulationId = signal<string | null>(null);
  private rsrpThreshold      = signal<number>(-140);
  readonly loading           = signal<boolean>(false);
  readonly error             = signal<string | null>(null);

  readonly filteredCoveragePoints = computed(() => {
    const sim = this.currentSimulation();
    const thr = this.rsrpThreshold();
    if (!sim) return [];
    return sim.coverage_points.filter(p => p.rsrp_dbm >= thr);
  });

  async loadSimulations(): Promise<void> {
    this.loading.set(true);  this.error.set(null);
    try {
      const summaries = await firstValueFrom(
        this.http.get<ApiSimulationSummary[]>(API_BASE));
      const partial: Simulation[] = summaries.map(s => ({
        simulation_id: s.id, metadata: s.metadata,
        coverage_area: { width_km: 3, height_km: 3, grid_spacing_m: 50 },
        coverage_points: []
      }));
      this.simulations.set(partial);
      if (partial.length > 0)
        await this.selectSimulation(partial[0].simulation_id);
    } catch (e) {
      this.error.set('No se pudo conectar con el backend.');
    } finally { this.loading.set(false); }
  }
}
\end{lstlisting}

El flujo de carga opera en dos fases: primero se solicitan los resúmenes (sin
\texttt{coveragePoints}) para rellenar el selector del \textit{sidebar}; al
seleccionar una simulación, se pide el JSON completo con los puntos de cobertura
y se actualiza la lista reactiva con \texttt{simulations.update()}. Este patrón
evita transferir datos pesados hasta que el usuario los necesita realmente.

La función \texttt{normalizeRsrp} convierte el valor RSRP al rango $[0, 1]$ para
el gradiente del \textit{heatmap} mediante la expresión
$(\text{RSRP} + 140) / 100$, donde $-140$~dBm se mapea a 0 y $-40$~dBm a 1.

\subsection{Componente de mapa}
\label{sec:frontend_mapa}

\texttt{MapComponent} integra Leaflet~1.9.4 con el \textit{plugin}
\texttt{leaflet.heat} para renderizar el mapa de calor. El usuario puede elegir entre
cartografía OpenStreetMap, imágenes de ESRI y relieve de OpenTopoMap.

El mapa se inicializa en \texttt{ngOnInit} centrado en las coordenadas del campus
URJC (40,2897°~N, 3,8244°~W) con zoom 13. Un \texttt{effect()} de Angular observa
los cambios en \texttt{filteredCoveragePoints} y en \texttt{currentSimulation}: cada
vez que alguno de esos \textit{signals} cambia, se reconstruye el \textit{heatmap}
automáticamente.

\begin{lstlisting}[language=JavaScript, caption={Renderizado del \textit{heatmap} de cobertura.}, label={cod:heatmap}]
private updateHeatmap(points: CoveragePoint[]): void {
  if (this.heatmapLayer) {
    this.map.removeLayer(this.heatmapLayer);
    this.heatmapLayer = null;
  }
  if (!this.heatmapVisible || points.length === 0) return;

  const heatData: [number, number, number][] = points.map(p => [
    p.latitude,
    p.longitude,
    this.simulationService.normalizeRsrp(p.rsrp_dbm)  // (rsrp+140)/100
  ]);

  this.heatmapLayer = (L as any).heatLayer(heatData, {
    radius: 18,  blur: 15,  maxZoom: 17,  max: 1.0,
    gradient: {
      0.0:  '#ef4444',   // Debil    (< -100 dBm)
      0.4:  '#f97316',   // Aceptable(-100..-85)
      0.55: '#fbbf24',   // Bueno    (-85..-70)
      0.7:  '#10b981'    // Excelente (> -70 dBm)
    }
  });
  this.heatmapLayer.addTo(this.map);
}
\end{lstlisting}

El gradiente de colores replica los umbrales de clasificación del
Cuadro~\ref{tab:rsrp_umbrales}: rojo para señal débil, naranja para aceptable,
amarillo para buena y verde para excelente. \texttt{leaflet.heat} dibuja todos los
puntos en un único elemento \texttt{<canvas>}. No se ha medido el tiempo de renderizado
del mapa en un navegador ni se afirma el cumplimiento de un umbral de rendimiento.

El manejador de \textit{click} del mapa busca el punto de cobertura más cercano
al lugar pulsado mediante distancia euclidiana en grados. Si la distancia es inferior
a 0,003° (aproximadamente 300~m), selecciona el punto y el \texttt{InfoPanelComponent}
muestra su RSRP, coordenadas y \textit{path loss} estimado. En caso contrario, se
deselecciona cualquier punto previo.

\subsection{Selector de cartografía y tiles vectoriales locales}
\label{sec:frontend_tile_themes}

La cartografía base condiciona la legibilidad del \textit{heatmap}: un fondo demasiado
saturado oculta los gradientes de cobertura, mientras que un fondo neutro los realza.
La aplicación incorpora un \texttt{MapTileService} que centraliza la gestión de las capas
cartográficas y un selector visual integrado en el mapa.

\subsubsection{Catálogo de cartografías}

El servicio define tres temas raster seleccionables desde la interfaz, todos ellos
servidos desde CDN públicos sin requerir clave de API:

\begin{table}[h!]
\centering
\caption{Cartografías disponibles en el frontend.}
\label{tab:cartografias}
\small
\begin{tabular}{|l|l|l|l|}
\hline
\textbf{Identificador} & \textbf{Proveedor} & \textbf{Tipo} & \textbf{Uso} \\
\hline
\texttt{osm}          & OpenStreetMap                & Raster PNG    & Cartografía estándar \\
\texttt{satellite}    & ESRI World Imagery           & Raster JPEG   & Vista aérea \\
\texttt{terrain}      & OpenTopoMap                  & Raster PNG    & Relieve y curvas de nivel \\
\hline
\end{tabular}
\end{table}

La elección persiste en \texttt{localStorage} bajo la clave \texttt{gd5g-map-theme},
de manera que sobrevive a recargas del navegador. Un \textit{signal} \texttt{current}
expone el tema activo y los componentes \texttt{MapComponent} y \texttt{ComparisonViewComponent}
reaccionan automáticamente reemplazando la capa de \textit{tiles} mediante un
\texttt{effect()} de Angular. Los tres proveedores raster requieren conectividad para
solicitar sus tiles; no se ha verificado un modo de funcionamiento sin conexión.

\subsubsection{Tiles vectoriales locales con PMTiles}

Las tres cartografías disponibles son \textit{tiles} raster servidos desde proveedores
externos. El servicio contiene una rama de código preparada para renderizar PMTiles, pero
el catálogo activo no define ningún tema vectorial; por ello, esta funcionalidad no se
ofrece desde la interfaz ni se considera una capacidad de uso implementada.

\paragraph{El formato PMTiles.}
PMTiles es un contenedor de \textit{tiles} desarrollado por Protomaps que almacena
millones de \textit{tiles} en un único fichero binario. El cliente accede a los
\textit{tiles} mediante peticiones HTTP \emph{range requests} (cabecera
\texttt{Range: bytes=offset-end}), por lo que no se requiere un servidor de
\textit{tiles} dedicado: basta con servir el fichero como contenido estático.
Su uso requeriría disponer del fichero y añadir un tema vectorial al catálogo; estas
condiciones no forman parte del flujo verificado en esta versión.

\paragraph{Renderizado vectorial sobre Leaflet.}
La dependencia \texttt{protomaps-leaflet} y la rama \texttt{kind: 'vector'} preparan la
integración con Leaflet, pero no hay una cartografía PMTiles declarada en
\texttt{MAP\_THEMES}. El fragmento siguiente ilustra esa rama de código, no una opción
activa del usuario:

\begin{lstlisting}[language=JavaScript, caption={Construcción polimórfica de la capa de \textit{tiles}.}, label={cod:tile_layer}]
createTileLayer(themeId?: string): L.Layer {
  const theme = MAP_THEMES.find(t => t.id === themeId) ?? this.current();

  if (theme.kind === 'raster') {
    return L.tileLayer(theme.url, theme.options);
  }

  // Vector: PMTiles via protomaps-leaflet
  return protomapsL.leafletLayer({
    url: this.resolvePmtilesUrl(theme.url),
    theme: theme.protomapsTheme ?? 'light'
  });
}
\end{lstlisting}

El código puede intentar usar primero el fichero local y recurrir a un CDN de Protomaps si
la comprobación HTTP falla. Como no existe un tema vectorial activo y no se ha verificado
el uso de un fichero local, este mecanismo se documenta como preparación técnica, no como
función disponible. Tampoco se han medido latencia, consumo de datos, tamaño de tiles o
funcionamiento sin conexión; no se atribuyen ventajas de rendimiento a esta rama.

\subsubsection{Mini-dashboard de selección}

La interfaz expone los tres temas raster mediante una pequeña tarjeta superpuesta
en la esquina inferior derecha del mapa, con un botón por cartografía representado
con un icono SVG distintivo. Al pasar el ratón se muestra el nombre y la descripción
detallada en un \textit{tooltip}, y el botón activo se resalta con el color de
acento del tema. En la vista comparativa de simulaciones (\texttt{ComparisonViewComponent})
cada uno de los dos mapas dispone de su propio selector independiente, lo que permite,
por ejemplo, contrastar la misma simulación sobre cartografía estándar de
OpenStreetMap y sobre la vista satélite de ESRI para evaluar la legibilidad del
\textit{heatmap} en función del fondo.

\subsection{Interfaz de usuario}
\label{sec:frontend_ui}

La aplicación ofrece un cambio de tema claro/oscuro gestionado por
\texttt{ThemeService}. Este servicio mantiene un \textit{signal} booleano
(\texttt{darkMode}), persiste la preferencia en \texttt{localStorage} con la clave
\texttt{gd5g-theme} y aplica la clase CSS \texttt{theme-light} al \texttt{<body>}
cuando el usuario alterna al modo claro. La cartografía activa se puede cambiar de forma
independiente mediante el selector del mapa.

El \textit{sidebar} ofrece un selector de simulaciones y un deslizador
(\textit{slider}) de rango $[-140, -40]$~dBm que filtra los puntos visibles en el
\textit{heatmap}. El formulario permite configurar los parámetros de simulación y
seleccionar en el mapa la ubicación del transmisor antes de enviar el
\texttt{POST /api/simulations}. La interfaz consulta la extensión del raster y avisa
cuando la ubicación o el área parecen rebasarla; esta advertencia es parcial, pues el
backend conserva la validación definitiva y puede rechazar perfiles fuera del GeoTIFF o
sin elevación válida.

El \texttt{InfoPanelComponent} se activa al hacer clic en un punto del mapa y muestra
el nivel RSRP con un indicador visual de barras (1 a 4 segmentos según el intervalo
visual del prototipo), las coordenadas geográficas exactas, los intervalos visuales
propios del prototipo y una estimación del
\textit{path loss}: $PL \approx \text{EIRP} - \text{RSRP}$, donde $\text{EIRP}$ se
aproxima como $P_{tx} + 13$~dB a partir de los metadatos de la simulación.

\subsection{Gestión del estado con \textit{signals}}
\label{sec:frontend_signals}

Angular introdujo los \textit{signals} en la versión~16 como alternativa a la gestión
de estado basada en \texttt{Observable} de RxJS o librerías externas como NgRx. En
esta aplicación hemos adoptado \textit{signals} de forma integral por tres razones:

\begin{enumerate}
\item \textbf{Simplicidad}: un \texttt{WritableSignal<T>} se lee como una función
  (\texttt{signal()}) y se escribe con \texttt{signal.set()} o
  \texttt{signal.update()}, sin necesidad de suscripciones ni operadores.
\item \textbf{Derivación reactiva}: \texttt{computed()} crea valores derivados que
  Angular recalcula automáticamente cuando sus dependencias cambian. En nuestro caso,
  \texttt{filteredCoveragePoints} depende de \texttt{currentSimulation} y de
  \texttt{rsrpThreshold}, y se actualiza sin intervención manual.
\item \textbf{Efectos secundarios controlados}: \texttt{effect()} ejecuta una función
  cada vez que los \textit{signals} leídos dentro de ella cambian. Lo usamos en
  \texttt{MapComponent} para reconstruir el \textit{heatmap} cuando cambian los
  datos filtrados.
\end{enumerate}

El modelo de datos del frontend se define en \texttt{simulation.model.ts} mediante
interfaces TypeScript: \texttt{Simulation}, \texttt{CoveragePoint},
\texttt{SimulationMetadata} y \texttt{SimulationStats}. La función
\texttt{apiToSimulation()} se encarga de mapear el formato del backend (con campos
\texttt{lat}, \texttt{lng}, \texttt{rsrp}) al formato interno (con campos
\texttt{latitude}, \texttt{longitude}, \texttt{rsrp\_dbm}), desacoplando así el
modelo de la API del modelo de la vista.

\section{Integración del sistema}
\label{sec:impl_integracion}

Los componentes forman un flujo local de simulación, consulta y visualización que no
requiere una base de datos ni una cola de mensajes. La persistencia en ficheros permite
consultar resultados, aunque no implementa historial versionado ni trazabilidad de medidas.

\subsection{Flujo de datos extremo a extremo}
\label{sec:integracion_flujo}

El sistema mantiene dos caminos de cálculo que convergen en la visualización:

\begin{enumerate}
\item \textbf{Referencia MATLAB}: al ejecutar \texttt{simulacion\_cobertura\_5g.m},
  el bloque de exportación escribe un JSON en \texttt{matlab/output/}. Los escenarios
  existentes usan Longley--Rice y declaran \texttt{SRTM\_90m\_global}.
\item \textbf{Predicción interactiva Java}: la interfaz envía los parámetros con
  \texttt{POST /api/simulations}; Spring Boot valida la petición, calcula FSPL+Deygout
  usando el GeoTIFF local y persiste la respuesta JSON en el mismo directorio configurado.
\item \textbf{Consulta}: Angular solicita el listado mediante
  \texttt{GET /api/simulations} y el detalle completo mediante
  \texttt{GET /api/simulations/\{uuid\}}.
\item \textbf{Visualización}: la aplicación transforma los \texttt{coveragePoints},
  presenta el mapa de calor en Leaflet y permite inspeccionar o comparar escenarios.
\end{enumerate}

Los dos motores comparten el formato de salida general, pero sus metadatos identifican
modelo y terreno; los resultados no se deben mezclar como si procedieran de un mismo motor.

\subsection{Formato de intercambio JSON}
\label{sec:integracion_json}

El fichero JSON constituye el contrato de intercambio entre los motores y Spring Boot.
Las salidas contienen los campos raíz siguientes:

\begin{table}[H]
\centering
\caption{Esquema del fichero JSON de simulación.}
\label{tab:json_schema}
\begin{tabularx}{\textwidth}{|l|l|X|}
\hline
\textbf{Campo} & \textbf{Tipo} & \textbf{Descripción} \\
\hline
\texttt{id}             & \texttt{string}  & UUID~v4 generado por MATLAB o Java \\
\texttt{name}           & \texttt{string}  & Nombre descriptivo del escenario \\
\texttt{createdAt}      & \texttt{string}  & Marca temporal ISO~8601 (UTC) \\
\texttt{metadata}       & \texttt{object}  & Parámetros normalizados (frecuencia, modelo, etc.) \\
\texttt{parameters}     & \texttt{object}  & Parámetros de la ejecución; los nombres varían entre motores \\
\texttt{coveragePoints} & \texttt{array}   & Array de \texttt{\{lat, lng, rsrp\}} \\
\hline
\end{tabularx}
\end{table}

El campo \texttt{metadata} expone los parámetros en claves normalizadas en inglés
(\texttt{frequency\_ghz}, \texttt{tx\_power\_dbm}, \texttt{antenna\_height\_m},
\texttt{tx\_location}, \texttt{model}, \texttt{terrain\_data}) para que el frontend
pueda consumirlos directamente. El campo \texttt{parameters} conserva el conjunto de
parámetros empleado por cada motor; sus nombres no son idénticos entre MATLAB y Java.

En el lado del backend, la deserialización la realiza Jackson a través de
\texttt{ObjectMapper.readTree()}, que construye un árbol \texttt{JsonNode} genérico.
Para el listado, \texttt{parseSummary} extrae solo \texttt{id}, \texttt{name},
\texttt{createdAt} y \texttt{metadata} del nodo raíz, evitando cargar en memoria el
array \texttt{coveragePoints} que puede contener miles de elementos.

En el lado del frontend, la función \texttt{apiToSimulation()} (definida en
\texttt{simulation.model.ts}) mapea los campos del JSON a la interfaz
\texttt{Simulation}, renombrando \texttt{lat}/\texttt{lng}/\texttt{rsrp} a
\texttt{latitude}/\texttt{longitude}/\texttt{rsrp\_dbm} para mantener la coherencia
con el modelo interno de la aplicación.

\subsection{Configuración del entorno de desarrollo}
\label{sec:integracion_entorno}

Para ejecutar el sistema completo en un entorno de desarrollo local se necesitan los
siguientes prerrequisitos:

\begin{table}[H]
\centering
\caption{Prerrequisitos del entorno de desarrollo.}
\label{tab:prereqs}
\begin{tabular}{|l|l|l|}
\hline
\textbf{Herramienta} & \textbf{Versión mínima} & \textbf{Propósito} \\
\hline
MATLAB          & R2025b  & Generación de escenarios de referencia \\
Java JDK        & 17      & Compilación y ejecución del backend \\
Apache Maven    & 3.9+    & Gestión de dependencias del backend \\
Node.js         & Según \texttt{package.json} & Entorno de ejecución del frontend \\
Angular CLI     & Según \texttt{package.json} & Compilación y servidor de desarrollo \\
\hline
\end{tabular}
\end{table}

El orden de arranque recomendado es:

\begin{enumerate}
\item Opcional: ejecutar la simulación MATLAB para generar un escenario de referencia.
  La simulación bajo demanda de la aplicación usa el motor Java y no necesita MATLAB.
\item Arrancar el backend desde la carpeta \texttt{backend/} con
  \texttt{mvn spring-boot:run}. El servidor queda disponible en
  \texttt{http://localhost:8080}.
\item Arrancar el frontend desde \texttt{frontend/} con \texttt{ng serve}. La
  aplicación queda disponible en \texttt{http://localhost:4200}.
\end{enumerate}

La configuración CORS del backend (sección~\ref{sec:backend_config}) permite que el
frontend en el puerto 4200 realice peticiones al backend en el puerto 8080 sin que
el navegador las bloquee. En un despliegue unificado (por ejemplo, sirviendo los
\textit{assets} estáticos de Angular desde el propio Spring Boot), esta configuración
dejaría de ser necesaria.

\subsection{Despliegue}
\label{sec:integracion_despliegue}

El backend se empaqueta como un JAR ejecutable estándar de Spring Boot:

\begin{lstlisting}[caption={Compilación y ejecución del backend.}, label={cod:build}]
cd backend/
mvn clean package
java -jar target/gemelo5g-0.1.0-SNAPSHOT.jar \
     --simulations.data-dir=/ruta/a/jsons
\end{lstlisting}

El comando \texttt{mvn package} compila el código, ejecuta los tests y produce un JAR con todas las dependencias
incluidas gracias al \textit{plugin} \texttt{spring-boot-maven-plugin}. La propiedad
\texttt{simulations.data-dir} se puede sobreescribir por línea de comandos para
apuntar a cualquier directorio que contenga los JSON de simulación.

Para el frontend, \texttt{ng build} genera los \textit{assets} estáticos en
\texttt{dist/}, que pueden servirse con cualquier servidor HTTP (Nginx, Apache, o
el propio Spring Boot copiándolos a \texttt{src/main/resources/static/}).

En el contexto de este TFG, ambos servicios se ejecutan en la misma máquina durante
el desarrollo y las pruebas. Un despliegue táctico real seguiría el mismo esquema en
un equipo portátil de campo, donde la latencia de red es irrelevante y la prioridad
es la autonomía del sistema sin dependencias de conectividad externa.
